\section{El impacto de la IA en la ingeniería de software: se necesita más Systems Judgment}

El curso cierra con la educación y el futuro de los IDE. La postura de Asif no es que «los fundamentos de CS hayan perdido vigencia». Sostiene que, a medida que aumenta el código generado, importa más la capacidad de entender state, performance, security, failure e interface. Un Agent puede escribir más código, pero también crea más rápido problemas de dependencias, reintentos y permisos. El trabajo del ingeniero ya no se limita a escribir línea por línea: también abarca definir el contract, verificar la evidence y gestionar el change.

\subsection{El IDE se convierte en un Agent workspace}

El IDE alojará varios agent, background task, terminal, browser, test y deployment feedback. La infraestructura debe permitir que el usuario sepa qué lee, qué modifica y qué ejecuta el agent. También debe conservar un control plane que permita pausar, revisar y revertir. Una autonomy mayor no debe costar la workspace state visibility.

\begin{importantbox}{Las cuatro capacidades centrales de la AI-native software engineering}
\begin{enumerate}
  \item Diseñar versioned interface, invariant y permission;
  \item establecer test, eval, diff y rollback para cada generated change;
  \item identificar los cuellos de botella del sistema en queue, database, network y provider;
  \item convertir los incident y el user feedback en guardrail, en lugar de solo corregir el prompt.
\end{enumerate}
\end{importantbox}

\subsection{AI-assisted incident response}

La IA puede resumir logs, relacionar deploy, generar query, explicar el runbook y proponer una hypothesis. Sin embargo, un incidente en producción ocurre con evidencia incompleta y con costos que cambian rápido. La IA no debe tener production authority ilimitada. El incident commander debe confirmar el objetivo, detener el feedback loop, elegir qué riesgo asumir y registrar las decisiones.

\lecturefigure{16-incident-learning.png}{La IA puede acelerar el procesamiento de la evidencia, pero las decisiones de producción siguen correspondiendo al human incident command.}{Subtítulos locales 47:50--48:37; diagrama conceptual redibujado.}

\begin{knowledgebox}{Lectura de la figura: qué paso acorta mejor el Assistant}
Puede encontrar correlaciones temporales en un volumen enorme de metrics/logs, generar una safe read-only query, comparar los deploy recientes u ordenar el historial del chat en una timeline. Los write de alto riesgo, el failover, la data repair y la customer communication siguen requiriendo un owner explícito y una approval.
\end{knowledgebox}

\begin{warningbox}{La causa raíz que genera la IA también requiere evidencia}
Un modelo de lenguaje tiende a convertir una correlación en un relato causal. Cada hypothesis debe enlazar a un metric, trace, log, config o reproduction. El postmortem debe registrar lo unknown y cada rejected hypothesis, para no tomar una explicación fluida por un hecho.
\end{warningbox}

\teachervoice{Al final de la clase, el docente plantea el AI-assisted incident response como una dirección posible. La conclusión más segura para llevar a otros contextos es esta: la IA aumenta la evidence bandwidth, y las personas conservan el objetivo, el riesgo y la authority.}

\subsection{Resumen de la sección}

El AI coding no le quita valor al systems knowledge. Al contrario: cuanto más capaz es la generación, más se vuelven centrales para el producto y para la carrera profesional la verificación, el aislamiento, la recuperación, la seguridad y la contabilidad de recursos.
